\documentclass[../main.tex]{subfiles}
\begin{document}
%==========================================TIÊU ĐỀ TẬP========================================
\begin{table}[h]
  \begin{adjustwidth}{-1cm}{}
    \begin{tabular}{>{\centering\arraybackslash}p{6cm}|>{\raggedright\arraybackslash}p{12.5cm}}
      \multirow{5}{6cm}
      % Cột bên trái: ÉPISODE và số tập
      {\\[-40pt]
        \flaregothic\fontsize{20pt}{20pt}\selectfont \centering
        {\contour{errortwo}{\textcolor{black}{ÉPISODE}}}\color{black}\\
        \burbank\fontsize{72pt}{72pt}\selectfont
        \includegraphics[height=3cm]{../../../../Image/Book?/number/num4.png}
        \\[18pt]
        % \flaregothic\fontsize{14pt}{14pt}\selectfont
        % \color{epattr}BIENVENUE, \uline{\color{Banpire} Mel} !
        % \flaregothic\fontsize{18pt}{18pt}\selectfont
        % \color{epattr}ÉPISODE PERDU
        \color{black}
      }
       &
      % Dòng thứ nhất: tên tập tiếng Nhật
      {\fotskip\fontsize{18pt}{18pt}\selectfont \ruby{帰}{かえ}りましょう\ruby{転}{てん}\ruby{校}{こう}\ruby{生}{せい}
      }  \\[6pt]
      % Dòng thứ hai: tên tập tiếng Anh
       & {\cafeteria\fontsize{18pt}{18pt}\selectfont Let''s Go Home, Transfer Student}                \\[4pt]
      % Dòng thứ ba: tên tập tiếng Việt
       & {\vnmsans\fontsize{18pt}{18pt}\selectfont Ta cùng về nhà thôi, những học sinh chuyển trường} \\[8pt]
      % Dòng thứ tư: Nhân vật xuất hiện
       & {\caviar{\fontsize{14pt}{14pt}\selectfont Track used: The Piano}}
      \\[6pt]
      % Dòng cuối cùng: Thời gian diễn ra của tập (thường sẽ là ngày phát hành)
       & {\continuum\fontsize{16pt}{16pt}\selectfont Ngày phát hành: 24/03/2022}                      \\[8pt]
    \end{tabular}
  \end{adjustwidth}
\end{table}
%===========================================NỘI DUNG==========================================
\color{black}

\pov{Haragen Saikyou}

``Vậy\dots\ hai người có vẻ như là các học sinh mới chuyển đến à.''

Shino nhìn chúng tôi thêm một lúc, rồi khóe môi cong lên. Không phải nụ cười chào hỏi. Không phải nụ cười vui vẻ. Mà là thứ gì đó vừa nham hiểm, vừa như đang thử thách.

Nhưng\dots\ ít nhất lúc này, tôi không cảm nhận được ý định hại người từ cô ấy. Không như trong giấc mơ, như cách mà ả đã làm với Nii-chan của tôi.

Cô ta vẫn đứng đó, lặng lẽ quan sát. Ánh mắt xanh của cô như đang dò tìm điều gì đó bên trong chúng tôi, thứ không thể che giấu dù có cố gắng. Tôi và Nii-chan bất giác giữ im lặng, chẳng biết nên lên tiếng hay chờ cô ấy chủ động.

Vài giây\dots\ hoặc có thể là lâu hơn thế, Shino bỗng quay đầu lại, ánh mắt hướng về khoảng hành lang trống trải phía trước.

``N-Này! Đợi đã!''

Không một lời, cô bắt đầu bước đi, chậm rãi nhưng dứt khoát, bóng lưng khuất dần sau góc rẽ, bỏ mặc những tiếng gọi của tôi.

Chúng tôi vẫn đứng yên, lắng nghe tiếng bước chân ấy nhỏ lại cho tới khi biến mất hoàn toàn.

``\textit{\dots Tại sao Shino lại ở đây? Mình tưởng là ả ta chỉ xuất hiện trong mơ thôi chứ\dots?}''  câu hỏi bật ra trong đầu tôi, nhưng không ai trả lời. Hoặc đúng hơn, không ai biết được câu trả lời, và cũng không có ai ở đây ngoài chúng tôi trả lời.

Nii-chan cũng đang nghĩ điều tương tự, tôi có thể đoán ra từ nét mặt anh. Và cùng với câu hỏi đó, một mối nghi ngờ khác len vào, nặng nề như cơn sương mù trước bình minh, và cũng là thứ mà Saki-san thắc mắc: ``\textit{Lời nguyền đó\dots\ thực sự bắt đầu từ đâu?}''

Không gian im lặng quanh chúng tôi bỗng trở nên đặc quánh. Chỉ còn lại tiếng thở của hai anh em, và cảm giác rằng hành lang này vừa nuốt mất những câu trả lời mà đáng lẽ chúng tôi phải tìm được.

Tôi còn chưa kịp hoàn hồn với sự xuất hiện của Shino, chỉ để rồi cô ta lẳng lặng bỏ đi, không để chúng tôi kịp hỏi một câu nào\dots

*XOẠCH!*

Một tiếng cửa trượt vang lên ở ngay bên cạnh, chói tai đến mức làm tôi rùng mình. Cả tôi và Nii-chan cùng giật bắn người. Từ bên trong lớp 1-C, vốn dĩ đã tan hoang thành đống đất đá, một bóng người chậm rãi bước ra.

\img{e1-4a}

Hanabi-san.

Cái tên bật ra trong đầu tôi ngay tức khắc. Nhưng\dots\ có gì đó sai lạ. Rõ ràng cô ấy chính là Hanabi-san, thế nhưng\dots\ tại sao sau trận động đất kinh hoàng ấy, cô lại bước ra với dáng vẻ hoàn toàn nguyên vẹn?

Bộ đồng phục trắng muốt không một vết bụi, mái tóc vẫn suôn mượt như vừa được chải gọn gàng, làn da không có lấy một vết xước nhỏ.

``\textbf{H-Hanabi-san!}'' tôi và Nii-chan đồng thanh, ngạc nhiên tới mức gần như nghẹn lời. Trái tim tôi đập dồn dập, như thể có hy vọng nào đó bùng lên trong phút chốc. Tôi lao tới, mặc kệ tiếng gọi với theo của Nii-chan.

``\textbf{N-Này, Saikyou! Khoan đã!}''

Tôi sải bước đến gần, định ôm chầm lấy Hanabi-san trong niềm vui sướng vì cô ấy vẫn bình an. Nhưng rồi\dots\ tôi chợt khựng lại, kiềm chế bản thân, giữ một khoảng cách nhỏ.

``\dots Hanabi-san?'' tôi lên tiếng, giọng hơi run.

Cái gì đó không ổn. Tôi cảm thấy rõ điều đó ngay khi ánh mắt tôi chạm vào đôi mắt của cô ấy.

Không phải ánh mắt lo lắng, cũng chẳng phải vui mừng khi thấy chúng tôi. Đôi mắt của cô trống rỗng, vô hồn, lạnh lẽo đến mức khiến tôi cảm giác như mình vừa chạm phải một bức tường băng.

Và rồi, Hanabi-san khẽ mỉm cười. Một nụ cười máy móc.

``Tên mình là Hanabi! Còn đây là Inari!''

\dots Cái gì cơ?

Tôi đứng sững, tim hẫng một nhịp.

``C-Cậu đang nói gì thế, Hanabi-san? Bọn mình quen nhau rồi mà! Hay\dots\ hay là cơn động đất làm cậu\dots\ chập mạch luôn rồi!?'' tôi lắp bắp, cố nặn ra một lời giải thích hợp lý nhất. Nhưng càng nhìn cô ấy, mồ hôi lạnh càng túa ra trên lưng tôi.

Inari-san vốn dĩ luôn đứng cạnh cô ấy, như thể họ là cặp bài trùng. Nhưng Hanabi-san lại tự giới thiệu cả hai, như thể chúng tôi chưa từng gặp mặt.

Và rồi, thay vì trả lời, Hanabi-san lại mở miệng, lặp lại y nguyên từng chữ, từng nhịp điệu: ``Tên mình là Hanabi! Còn đây là Inari!''

Lần thứ hai.

Lần thứ ba.

Giọng nói đều đều, như một con búp bê được cài sẵn một đoạn băng cassette.

``\dots Chết tiệt\dots'' Nii-chan khẽ buông ra, giọng đầy chua chát nhưng cũng không kém phần hãi hùng: ``Chẳng khác gì mấy con NPC trong game cả\dots\ chỉ biết lặp đi lặp lại một câu thoại duy nhất.''

Tôi quay sang nhìn anh, sửng sốt. Cách anh nói\dots\ nghe thì có lý, nhưng chính vì thế mà càng khiến máu trong người tôi lạnh đi. Hanabi-san\dots\ không còn là Hanabi-san mà tôi biết.

Bất giác, toàn thân tôi run lên. ``Nii-chan\dots\ đừng có nói thế chứ\dots'' Tôi lùi lại một bước, hoảng sợ như thể chính ý nghĩ ấy đã biến nó thành sự thật.

Nhưng trước khi tôi kịp nói thêm gì, Nii-chan bất ngờ thốt lớn, tay chỉ ra phía sau tôi: ``\textbf{Đ-Đằng sau em kìa, Saikyou!}''

Tôi giật mình quay ngoắt đầu lại.

Một bóng người đã đứng đó từ bao giờ.

\img{e1-4b}

Inari-san.

Cổ họng tôi ứ lại, nhìn thấy một người bạn quen thuộc nữa. ``I-Inari-san! C-Cậu--''

``Khoan đã, Saikyou!'' Nii-chan kịp ngăn tôi lại, giọng gấp gáp. ``Cô ấy\dots\ trông cũng không ổn đâu.''

Tôi nín thở, nhìn kỹ hơn. Và rồi\dots\ trái tim tôi thắt chặt.

``\dots Trời đất ơi\dots\ Cả cậu cũng thế sao, Inari-san\dots!?''

Đôi mắt ấy. Đúng, lại là đôi mắt ấy. Trống rỗng, vô hồn, giống hệt như Hanabi-san. Cơ thể hoàn toàn nguyên vẹn, không một vết thương, không một dấu hiệu vừa trải qua thảm họa nào. Nhưng linh hồn bên trong\dots\ như đã biến mất hoàn toàn. Trông cô ấy không khác gì như một con búp bê - như lời mà Touka-san đã khen tôi ngày đầu tiên tôi đến trường.

Chỉ khác rằng\dots\ con búp bê mà cô ấy nói\dots\ giờ đã theo đúng nghĩa đen. Không có sức sống, chỉ còn vỏ ngoài bóng bẩy.

Inari-san mở miệng.

``Học sinh mới chuyển tới phải được đối xử thật tốt\dots''

Từng từ rơi xuống, nặng nề, đơn điệu, giống như một chiếc kim gõ liên hồi vào thái dương tôi.

``Nếu không\dots\ cơn đại họa sẽ ập xuống\dots''

Và cũng giống như cô bạn của mình, cô bạn tóc trắng ấy cũng lặp đi lặp lại câu nói ấy.

Tôi chết sững.

Không\dots\ không thể nào\dots

Đó chẳng phải chính là ``lời đồn'' sao? Câu nói mà tôi đã nghe trong giấc mơ, đọc trong những mẩu báo vàng ố? Câu mà Inari-san đã thốt lên trong hiện thực, khi trận động đất khởi phát?

Bàn tay tôi run rẩy bấu chặt lấy tay áo mình.

Tất cả những gì vừa xảy ra\dots\ có phải\dots\ đều bắt nguồn từ cú đập lưng mạnh tay của Suzune-san hôm ấy? Lúc mà Nii-chan khẽ nhăn mặt, và mọi người nghĩ rằng cậu khó chịu? Chẳng lẽ chỉ vì một việc nhỏ bé đó\dots\ mà họ bị coi là đã ``không đối xử tốt'' với một học sinh mới chuyển đến?

Nỗi bất an, khiếp đảm, dồn lại trong lồng ngực. Tôi ngước nhìn Nii-chan, tìm chút hy vọng từ ánh mắt anh. Nhưng thứ tôi thấy chỉ là một sự nghi hoặc căng thẳng\dots\ và nỗi sợ cũng chẳng khác tôi là bao.

``Lời nguyền\dots\ thật sự là thật sao\dots?'' tôi thì thầm, môi run lên, không biết mình đang hỏi anh hay hỏi chính bản thân mình.

Hanabi-san và Inari-san - những người bạn vốn dĩ thân thuộc - nay bỗng nhiên như biến thành thứ gì đó xa lạ. Họ đứng lặng, mắt vô hồn, miệng lặp lại cùng một cụm từ như thể bị ai đó giật dây, khiến cho hành lang bỗng chốc trở nên lạnh lẽo và méo mó lạ thường. Tiếng vang dội trong không gian hẹp biến mọi lời nói thành một bản nhạc ma mị, khiến Saikyou dựng hết cả tóc gáy.

Không khí trong hành lang lúc này đặc quánh lại, như thể từng phân tử không khí đều bị đông cứng bởi sự lặp lại vô hồn của hai người bạn. Mỗi lần cô bạn Nhiệt huyết cất tiếng, câu giới thiệu máy móc vang lên, âm thanh vang vọng giữa những bức tường trống trải, tạo thành một nhịp điệu rối loạn, bất thường. Còn Lớp trưởng thì vẫn đều đều nhắc lại ``lời nguyền'', giọng nói của cô lạnh lẽo, không mang chút cảm xúc nào, như thể đó là một đoạn ghi âm cũ kỹ được phát lại không ngừng.

Tiếng nói của họ hòa quyện vào nhau, tạo thành một bản hợp xướng kỳ dị, khiến cho hành lang vốn đã vắng lặng càng trở nên ngột ngạt, u ám. Ánh sáng từ cửa sổ hắt vào chỉ càng làm nổi bật sự bất động của hai bóng người, khiến mọi thứ xung quanh như bị bóp méo, xa lạ. Tôi cảm giác như mình đang đứng giữa ranh giới của thực và mơ, nơi mọi quy luật bình thường đều bị đảo lộn bởi sự hiện diện của cái gọi là lời đồn ở ngôi trường này.

Tôi nhìn hai cô bạn, môi họ vẫn lặp lại những câu nói vô hồn, ánh mắt trống rỗng như hai con rối bị giật dây. Tim tôi co thắt lại. Đó\dots\ không còn là họ nữa. Những người bạn từng cười nói với tôi, giờ chỉ còn là vỏ rỗng lập đi lập lại một âm thanh lạnh sống lưng. Tôi muốn lao đến, muốn lay họ dậy, muốn hét thật to để họ tỉnh lại\dots\ nhưng đôi chân tôi như bị ghì xuống, chỉ có thể bất lực nhìn họ.

``\textbf{Ta phải rời khỏi đây!}'' tiếng Nii-san gầm lên, giọng vang rền trong sự căng thẳng tột độ.

Nước mắt dâng lên nơi khóe mắt nhưng tôi không thể khóc. Tôi biết anh ấy đúng. Tôi biết tôi chẳng thể cứu nổi họ. Nhưng quay lưng bỏ lại hai người bạn, để mặc họ ở lại trong hình hài méo mó đó\dots\ tim tôi đau nhói, như thể bị xé làm đôi. Tôi nắm chặt tay, rồi buông ra, cuối cùng cắn răng chạy theo anh.

``Tôi xin lỗi hai người\dots\ tôi không thể cứu hai người được\dots''

Trong nỗi bứt rứt xé lòng, tôi buộc phải quay lưng, chạy theo ông anh trai, bỏ lại tất cả phía sau như một vết dao khắc sâu trong ký ức tốt đẹp sớm nở chóng tàn.

\ct

Chưa kịp lấy hơi, khi hai anh em chúng tôi vừa bắt đầu chạy, một âm thanh vang lên. Không phải thứ tiếng đều đều, vô hồn như của hai cô bạn đã ``chết'' kia, cũng chẳng ma mị đến lạnh người. Nó mang chất ``người'', đầy rung động, khiến tôi bất giác nhớ tới giọng nói của cô gái tóc nâu khi nãy.

Âm thanh phát ra từ phía cầu thang, ngay trước cánh cửa nặng nề kia. Nii-san nghiến răng, kéo cả tôi rẽ vào, nhưng chưa kịp phản ứng thì cả hai đã bị một lực mạnh kéo giật sang một bên.

``\textbf{Á!}''

Lực đó mạnh tới mức kéo tôi ngã nhào vào lòng đất bụi bặm. Anh tôi lập tức chống tay, định vùng dậy phản kháng\dots

``A-Ai đó!?''

\dots nhưng rồi đôi mắt anh mở to khi nhìn rõ kẻ vừa lôi họ: một cô gái có mái tóc nâu mềm mại, khuôn mặt y hệt Shino, nhưng đôi mắt lại là màu nâu. Và cũng mặc bộ đồng phục giống như tôi nữa, chỉ khác là đang khoác chiếc áo choàng đen bên ngoài.

Một trong hai người đã đưa chúng tôi thoát khỏi cơn ác mộng trong đêm hôm đó.

``S-Sora!? Tại sao cậu lại--'' tôi hoảng hốt, ngạc nhiên đến mức lời nói nghẹn lại.

Sora không để cô nói hết. Cô ấy dứt khoát cắt ngang, giọng hối hả: ``Tớ không còn thời gian giải thích cho hai cậu đâu! Mau chạy thôi!''

Ngay khi câu nói vừa dứt, mặt đất lại rung lắc dữ dội. Cả ngôi trường như gầm lên một lần nữa, tường trần rạn nứt, từng mảng gạch vụn rơi xuống như mưa. Chúng tôi lập tức hiểu rằng đây hẳn là dư chấn của trận động đất ban nãy, nhưng mức độ thì khủng khiếp hơn nhiều.

\img{e1-4c}

Thế là cả ba lao xuống cầu thang. Mỗi bước chạy nặng như đá, vừa sợ vừa khẩn trương. Hơi thở tôi dồn dập, tim đập loạn trong lồng ngực, mồ hôi lạnh rịn ra sau gáy. Tiếng nứt vỡ, tiếng gạch đá rơi vỡ tan xen lẫn trong tiếng chân đập thình thịch. Tôi không dám ngoái lại, không dám nhìn những gì mình đã bỏ lại phía sau. Chỉ còn một ý nghĩ duy nhất trong đầu: phải thoát ra khỏi nơi này, bằng mọi giá.

Từ phía trên, trần nhà bắt đầu nứt nẻ, những bức tường bắt đầu rơi lả tả, cùng với đó, những nắm đất rơi ào ào xuống như mưa trút nước. Những ô cửa kính ở hành lang thì toác ra, rơi thẳng xuống các tầng dưới.

Khung cảnh này\dots\ giống y như những gì diễn ra trong cơn ác mộng đêm hôm đó. Chỉ khác rằng, tôi cảm tưởng cơn động đất, hoặc dư chấn, còn mạnh hơn cả trong đó. Mọi tiếng động phía sau dần lùi xa, nhưng lại bị thay thế bằng những thứ còn đáng sợ hơn: tiếng cười khe khẽ, lanh lảnh, ngắt quãng\dots\ y hệt như trong cơn ác mộng đêm đó.

``Nii-chan\dots'' tôi thì thào, giọng run rẩy, ``có phải anh cũng nghe thấy không\dots?''

Nii-chan không trả lời ngay. Anh chỉ siết chặt bàn tay tôi hơn, buộc tôi phải tiếp tục chạy. Tôi nghe rõ tiếng tim mình đập dồn dập, xen lẫn cả tiếng tim anh. Rồi anh khẽ gằn giọng: ``Ừ\dots\ giống hệt như giấc mơ đó. Y chang\dots''

Tôi cắn môi đến bật máu. Tại sao mọi thứ lại trùng hợp đến mức này? Giấc mơ\dots\ hay đúng hơn, ác mộng\dots\ đang nhấn chìm chúng tôi trong hiện thực. Tôi quay đầu thoáng nhìn Sora đang bám sát phía sau, gương mặt cậu ấy bình tĩnh một cách lạ lùng giữa màn bụi mù.

Như thể đây không phải là lần đầu cô ấy trải qua chuyện này.

``Cứ chạy đi,'' cô khẽ nói, đủ để chúng tôi nghe thấy. ``Đừng dừng lại. Nếu dừng, các cậu sẽ bị nó nuốt chửng đấy.''

``Nó\dots?'' Tôi thốt lên, không kịp hỏi tiếp.

Không cần hỏi thêm nữa. Tôi và ông anh trai cũng mập mờ hiểu được \emph{nó} là cái gì.

Và rồi, như để trả lời, dọc hành lang vừa hiện ra trước mắt, tôi thấy chúng\dots\ những cái bóng.

Thoáng đầu tiên, một dáng người đứng nghiêng, tóc dài xõa, trên tóc có cài một bông hoa\dots Touka-san? Tôi rùng mình, nhắm chặt mắt rồi mở ra, nhưng bóng đó đã biến mất.

Tiếp theo, từ khung cửa lớp bên trái, lóe lên hình ảnh một cô gái với mái tóc hai búi tròn, là Suzune-san. Tiếng cười rít lên the thé, rồi tan biến như chưa từng có.

Trái tim tôi thắt lại. Từng khuôn mặt quen thuộc, từng người bạn đã gục ngã dưới lớp trần sập xuống\dots\ tại sao lại xuất hiện ở đây, như những hồn ma vương vấn?

``Không thể nào\dots'' tôi lẩm bẩm, đôi chân khựng lại nửa giây. Nhưng anh trai kéo giật tôi đi, miệng anh run lên: ``Cứ chạy đi, Saikyou! Đừng để ý tới họ!''

Tôi quay đầu. Một lớp học khác, ngay gần cầu thang, thấp thoáng một chiếc bờm tai cáo, Inari-san. Đằng xa, nơi ánh sáng bị bụi mờ che khuất, có mái tóc đuôi ngựa buộc lệnh, trông như Hanabi-san đang đứng nhìn chúng tôi. Và cuối cùng, ngay cạnh lan can tầng hai, một cái bóng nhỏ bé với hai bím tóc thấp, đung đưa trong gió bụi\dots\ không ai khác ngoài Saki-san.

Mồ hôi lạnh chảy dọc sống lưng. Tôi muốn hét lên, muốn hỏi họ tại sao lại xuất hiện như thế, nhưng cổ họng ứ lại. Những bóng hình đó không đáp lại, chỉ thoắt ẩn thoắt hiện, tựa như đang trêu ngươi chúng tôi bằng sự im lặng lạ lùng.

``\textbf{Saikyou-san! Kyuen-san! Tập trung vào! Đừng nhìn nữa!}'' tiếng Sora dõng dạc vang lên phía sau, không còn bình tĩnh như trước mà đã gấp gáp.

Chúng tôi gật đầu, chỉ biết cắm đầu chạy. Hành lang kéo dài vô tận, những cái bóng chập chờn bủa vây hai bên. Mỗi bước chân như nặng hơn gấp đôi, nhưng chẳng có lựa chọn nào khác.

Cuối cùng, hơi thở đã muốn đứt đoạn, cả ba chúng tôi cũng chạm được tới cầu thang. Tôi lao xuống, bước chân gấp gáp dội vang, khói bụi mờ ảo vây quanh. Tầng một đã ở ngay trước mặt, nhưng trong lòng tôi, nỗi sợ càng dâng cao.

Vì tất cả những gì đang xảy ra\dots\ thật sự không còn là ác mộng nữa.

\ct

Chúng tôi cắm đầu chạy. Tiếng trần nhà nứt rạn, những mảnh gạch vữa rơi lộp bộp xung quanh. Tôi vừa chạy vừa ngoái nhìn Nii-chan, lòng chỉ cầu mong anh ấy không bị gì. Ngôi trường này\dots\ cứ như đang tự nuốt chửng lấy chúng tôi.

``\textbf{Ra lối chính!}'' Nii-chan hô, kéo tay tôi chạy thẳng về phía sảnh lớn. Tôi gật đầu vội, tim đập dồn dập. Cửa kìa, ngay trước mắt\dots\ chỉ cần mở ra là thoát!

Chúng tôi lao đến, cùng nhau đẩy mạnh cánh cửa kính dày cộp ấy. Nhưng\dots\ nó không nhúc nhích. Tôi dồn hết sức, vai tê rần, Nii-chan cũng gồng hết lực, vậy mà cửa vẫn đóng chặt, như bị một thế lực nào đó khóa lại.

``Không\dots\ không thể nào\dots'' tôi thở hổn hển, hoảng loạn.

``Có ai đó khóa cánh cửa này sao!?'' anh tôi đập mạnh vào thành cửa, nhưng hoàn toàn không xi nhê gì.

``M-Mọi người! Vẫn còn một cửa thoát nữa!'' Sora ở bên cạnh tôi thốt lên, mồ hôi lấm tấm trên trán.

Nhưng ngay sau đó, một giọng nói khác giống như Sora chêm vào, chỉ có điều giọng ấy lại trầm và đơn điệu hơn, khiến máu tôi lạnh buốt.

\img{e1-4d}

``Thấy trường học có vui không? Dám cá là thấy vui đấy. Nhưng mà, tới giờ đi về rồi.''

Tôi quay phắt lại, tim như ngừng đập. Một cô gái tóc nâu, cài kẹp đỏ, đôi mắt xanh lạnh lẽo\dots\ đứng ngay giữa hành lang vỡ nát. Chính là cô gái mà tôi thoáng thấy trong lớp học trước khi động đất xảy ra.

Shino. Giọng cô ta trống rỗng như thể đang mơ, nghe mà rợn cả người. Tôi cảm thấy chân mình run lên, những ký ức về giấc mơ kinh hoàng ấy ùa về như thác lũ.

Nii-chan liền quát, giọng run rẩy nhưng cố tỏ ra mạnh mẽ: ``\textbf{Rốt cuộc thì chuyện gì đã xảy ra vậy!?}''

Tôi cũng không kìm được, hét theo, cổ họng nghẹn đắng: ``\textbf{Tại sao cô lại ở đây!? Tại sao lại giống y như trong giấc mơ!?}''

Nhưng cô gái ấy chẳng thèm trả lời. Chỉ mỉm cười nhạt, đôi môi mấp máy chậm rãi: ``Muộn rồi đấy. Nào, ta cùng đi về thôi. Nhanh lên.''

Tôi sững lại. Cô ta không nhìn tôi, cũng chẳng nhìn Nii-chan, mà ánh mắt lại hướng thẳng về phía người đang đi cùng chúng tôi, Sora. Tôi thấy rõ bàn tay Nii-chan siết chặt lại, như thể anh đang cố kìm nén nỗi sợ hãi đang dâng trào.

``Cậu\dots\ đang gọi tôi ư?'' Sora khẽ hỏi, giọng run nhưng vẫn cứng rắn.

Shino gật nhẹ, rồi thở dài: ``Cậu vẫn nhất quyết cứng đầu vậy sao? Đến giờ phút này mà vẫn muốn bám lấy những người không liên quan?''

Tôi thấy Sora nghiến răng, mắt lóe lên thứ ánh sáng kiên định: ``Họ không phải `không liên quan''! Tôi sẽ không bỏ rơi bạn bè của mình. Cho dù chỉ mới quen biết, nhưng\dots\ họ là những người bạn mới của tôi!''

Lời nói đó làm tôi sững lại. Trong đầu bỗng thoáng hiện cảnh trong mơ, cô gái ấy chìa tay cầu xin được làm bạn, và tôi cùng Nii-chan đã gật đầu chấp thuận, dù vẫn còn rất lưỡng lự.

Shino im lặng, rồi quay lưng bỏ đi, mặc cho trần nhà tiếp tục đổ sập từng mảng. Tôi muốn gọi với theo, muốn hỏi cho rõ mọi chuyện, nhưng Sora cắt ngang bằng cái lắc đầu, như muốn che giấu điều gì đó mà có lẽ chỉ hai người họ mới hiểu.

Ngay khoảnh khắc ấy, từ mặt sàn nứt toác phía trước, một cánh tay đen ngòm thò ra, nắm chặt lấy chân Nii-chan.

\img{e1-4e}

``!?'' anh kêu lên, loạng choạng mất thăng bằng.

Tôi chết sững, không kịp phản ứng gì. Một cánh tay nữa chụp lấy chân còn lại, rồi cả hai cùng kéo anh ấy xuống. Xuống cái hố tối om ấy.

``\textbf{Nii-chan!!!}'' tôi hét lạc giọng, lao tới nắm lấy tay anh trai. Không thể để mất anh ấy được. Không thể để anh ấy biến mất như những người bạn khác.

Nhưng rồi, Sora ngay lập tức túm lấy tay tôi, hoảng hốt: ``\textbf{Đừng! Cậu sẽ bị kéo theo đấy! Tớ không muốn mất thêm một người bạn nào nữa!}''

Trong chốc lát, nước mắt trào ra. Tôi ngoái lại cô ấy, đáp lại: ``Bạn bè có thể mất\dots\ nhưng Nii-chan thì không thể mất được!! Xin lỗi, Sora\dots\ nhưng tôi phải đi!''

Nói rồi, mặc cho cô ấy gào gọi, tôi lao thẳng xuống theo anh trai. Khoảnh khắc bóng tối nuốt chửng lấy tôi, tôi vẫn nghe rõ tiếng tim mình đập dồn, cùng tiếng hét xé lòng của Sora bên trên\dots

``\textbf{Đừng mà!!!}''

Và rồi, cô ấy cũng nhảy xuống. Tôi nghe thấy tiếng gió rít bên tai khi Sora lao xuống theo chúng tôi, áo choàng đen tung bay trong không trung.

``\textbf{Đừng có lo cho chúng tôi chứ, đồ ngốc!}'' tôi hét lên về phía cô nàng có đôi mắt nâu đó, giọng nghẹn ngào.

``Tớ không thể cứ thế mà để mặc các cậu được\dots'' cô ấy đáp lại, giọng run run nhưng kiên quyết. ``Các cậu là bạn của tớ. Và tớ đã thề sẽ không bao giờ bỏ rơi bạn bè nữa!''

Dù biết rõ bên dưới kia là thứ gì, cô vẫn không bỏ mặc chúng tôi. Tôi cảm thấy ấm áp lạ thường, dù đang rơi trong bóng tối vô định.

Ánh sáng ở trên lịm dần, lịm dần, rồi tắt hẳn.

Ba chúng tôi cùng rơi xuống, không biết điều gì đang chờ đợi phía dưới, nhưng ít ra\dots\ chúng tôi không hề đơn độc.

Trong khoảnh khắc đó, tôi bỗng nhận ra: có lẽ từ nay, mọi thứ sẽ không còn trở lại như trước nữa.

Và bóng tối nuốt chửng tất cả.

Chờ đợi ở dưới lỗ đen là một thứ gì đó\dots\ mà có lẽ sẽ khiến tôi sởn gai ốc.

\end{document}